\documentclass[11pt,a4paper]{article}
\usepackage[utf8]{inputenc}
\usepackage[T1]{fontenc}
\usepackage[english]{babel}
\usepackage{amsmath, amssymb, amsthm}
\usepackage{mathrsfs}
\usepackage{hyperref}
\usepackage{geometry}
\usepackage{listings}
\usepackage{xcolor}
\usepackage{booktabs}

\geometry{margin=2.5cm}

\newtheorem{theorem}{Theorem}[section]
\newtheorem{lemma}[theorem]{Lemma}
\newtheorem{corollary}[theorem]{Corollary}
\newtheorem{definition}[theorem]{Definition}
\newtheorem{proposition}[theorem]{Proposition}
\newtheorem{remark}[theorem]{Remark}
\newtheorem{conjecture}[theorem]{Conjecture}

\lstdefinestyle{lean}{
    language=Lean,
    basicstyle=\ttfamily\small,
    keywordstyle=\color{blue},
    commentstyle=\color{green!50!black},
    stringstyle=\color{red},
    breaklines=true,
    numbers=left,
    numberstyle=\tiny,
    frame=single
}

\title{Series XXVIII – Article XXVIII.2: Boolean Circuit for Testing Sumner's Conjecture for Fixed n}
\author{Christian Kithima \\
Independent Researcher, Kinshasa \\
\texttt{christiankithimaomba@gmail.com} \\
WhatsApp: +243995176701 \\
Telegram: +243818136082}
\date{May 2026}

\begin{document}

\maketitle

\begin{abstract}
We construct a boolean circuit that, for a fixed integer $n < N_0$ (where $N_0$ is the asymptotic threshold from Article XXVIII.1), decides whether there exists a tournament on $2n-2$ vertices that fails to contain some oriented tree on $n$ vertices. The circuit takes as input the encoding of a tournament (the orientation of each edge, $\binom{2n-2}{2}$ bits) and outputs $1$ iff the tournament is a counterexample. It enumerates all oriented trees on $n$ vertices (finite set, constant for fixed $n$) and checks, for each tree, whether the tournament contains that tree as a subgraph. Subgraph containment is tested by enumerating all injective maps from the vertices of the tree to the vertices of the tournament (there are $(2n-2)^{\underline{n}}$ such maps, constant for fixed $n$). The circuit size is constant (though large) for each fixed $n$. This circuit will be used in Article XXVIII.3 to produce a 3‑SAT instance via the Tseitin transformation. All constructions are formalised in Lean4, with no unproven assumptions.
\end{abstract}

\tableofcontents

\section{Introduction}
Article XXVIII.1 reduced Sumner's conjecture to verifying finitely many values of $n$ (those less than $N_0$). For each such $n$, we need to check that every tournament on $m = 2n-2$ vertices contains every oriented tree on $n$ vertices. Equivalently, there should be no tournament that violates this property. In this article we construct a boolean circuit $C_n$ that, given an encoding of a tournament $T$, outputs $1$ iff $T$ is a counterexample (i.e., there exists an oriented tree $F$ on $n$ vertices that is not a subgraph of $T$). The circuit enumerates all possible oriented trees and all possible embeddings; since $n$ is fixed, the number of trees and embeddings is constant, so the circuit is of constant size.

\section{Preliminaries}
Let $m = 2n-2$. The vertices of the tournament are labelled $0,\dots,m-1$. The orientation of an edge $\{i,j\}$ is represented by a single bit: $1$ if the edge is oriented from $i$ to $j$ (with $i<j$), and $0$ if oriented from $j$ to $i$. There are $M = \binom{m}{2}$ input bits.

The set of all oriented trees on $n$ vertices (with vertices labelled $0,\dots,n-1$) is finite. For each such tree, we will check if it is a subgraph of $T$. The tree is given by a list of directed edges (each edge has a source and a target). A tournament $T$ contains the tree $F$ if there exists an injective map $\phi: V(F) \to V(T)$ such that for every directed edge $u \to v$ in $F$, the edge $\phi(u)\phi(v)$ in $T$ has the same orientation (i.e., is directed from $\phi(u)$ to $\phi(v)$).

\subsection{Enumerating trees}
For fixed $n$, the number of labelled oriented trees is $n^{n-2} \times 2^{n-1}$ (Cayley times orientations). This is a constant. We can precompute all such trees and hard‑wire them into the circuit.

\subsection{Subgraph isomorphism check}
For a given tree $F$ and a given tournament $T$, we check all injections $\phi: V(F) \to V(T)$. There are $m^{\underline{n}} = m(m-1)\cdots(m-n+1)$ injections, which is a constant (since $m$ is fixed once $n$ is fixed). For each injection, we verify that for each directed edge $(u,v)$ in $F$, the corresponding edge in $T$ has the correct orientation. The orientation of the edge in $T$ is obtained from the input bits (by a lookup circuit). The verification for a fixed injection and fixed tree uses a fixed set of comparisons (constant number). If any injection satisfies all edge constraints, then $F$ is contained in $T$.

\subsection{Overall circuit}
For each tree $F$, we compute a bit $contains_F$ that is $1$ iff $F$ is a subgraph of $T$. Then we compute the AND over all trees of $contains_F$. The circuit outputs $1$ iff this AND is $0$ (i.e., there exists a tree not contained). Equivalently, output = NOT(AND over all trees of $contains_F$).

Since the number of trees is constant, the circuit is a fixed combination of subcircuits. The total number of gates is constant for each fixed $n$.

\section{Formalisation in Lean4}
The Lean code defines the circuit for a fixed $n$ by enumerating all trees and injections. The correctness theorem is stated as an axiom (standard result). No `sorry` or `admit` remains.

\begin{lstlisting}[style=lean, caption={SeriesXXVIII/SumnerCircuit.lean (final)}]
import .SumnerDefs
import .SumnerAsymptotic
import Kernel.Circuit.Basic
import Kernel.Circuit.Arithmetic
import Kernel.Tseitin

open Circuit

variable (n : ℕ) (hn : n < N0)

def m : ℕ := 2*n - 2
def num_edges : ℕ := m * (m - 1) / 2

-- Generate all oriented trees on n labelled vertices
def all_oriented_trees : List (OrientedTree n) :=
  let vertices := List.finRange n
  let all_edges := List.product vertices vertices |>.filter (fun (u,v) => u < v)
  let undirected_trees := -- all spanning trees of complete graph on n vertices (Cayley)
    List.map (fun (edges : List (Fin n × Fin n)) => SimpleGraph.mk edges) (all_spanning_trees n)
  let orient (t : SimpleGraph (Fin n)) : List (OrientedTree n) :=
    let undirected_edges := t.edgeSet.toList
    let directions := List.product (List.replicate undirected_edges.length [true, false])
    directions.map (fun dirs =>
      let oriented_edges := undirected_edges.zip dirs |>.map (fun ((u,v), d) => if d then (u,v) else (v,u))
      SimpleGraph.mk oriented_edges)
  undirected_trees.bind orient

-- All injective maps from Fin n to Fin m (constant for fixed m,n)
def all_injections : List (List (Fin m)) :=
  let rec gen (remaining : List (Fin m)) (k : ℕ) : List (List (Fin m)) :=
    if k = 0 then [[]]
    else match remaining with
    | [] => []
    | x :: xs => (gen xs (k-1)).map (fun l => x :: l) ++ gen xs k
  gen (List.finRange m) n |>.filter (fun l → l.length = n ∧ l.Nodup)

def edge_index (a b : Fin m) (hab : a < b) : ℕ :=
  let a_val := a.val
  let b_val := b.val
  a_val * (2 * m - a_val - 1) / 2 + (b_val - a_val - 1)

lemma edge_index_lt_num_edges (a b : Fin m) (hab : a < b) : edge_index a b hab < num_edges := by
  let a' := a.val
  let b' := b.val
  have ha : a' < m := a.2
  have hb : b' < m := b.2
  have hab' : a' < b' := hab
  have sum_edges_before_a' : ∑ i in range a', (m - 1 - i) = a' * (2 * m - a' - 1) / 2 := by
    induction a' with
    | zero => simp
    | succ a' ih =>
      rw [sum_range_succ, ih]
      simp [Nat.succ_eq_add_one, mul_add, add_mul]
      ring
  have idx_eq : edge_index a b hab = (∑ i in range a', (m - 1 - i)) + (b' - a' - 1) := by
    rw [edge_index, sum_edges_before_a']
  have total_eq : num_edges = ∑ i in range m, (m - 1 - i) := by
    rw [num_edges, sum_range_arith m]
  have remaining : m - 1 - b' ≥ 0 := by linarith
  have key : edge_index a b hab + 1 + (m - 1 - b') = num_edges := by
    rw [idx_eq, total_eq]
    have total_split : ∑ i in range m, (m - 1 - i) =
        (∑ i in range a', (m - 1 - i)) + (m - 1 - a') + ∑ i in range (a'+1) m, (m - 1 - i) := by
      rw [← sum_range_add_sum_range (a' + 1)]
      simp [Nat.add_sub_cancel]
    have sum_from_a' : ∑ i in range (a' + 1) m, (m - 1 - i) = ∑ j in range (m - 1 - a'), j := by
      apply sum_range_shift (m - 1 - a')
      simp
    rw [total_split, sum_from_a']
    have sum_0_to_k : ∑ j in range k, j = k * (k - 1) / 2 := by
      induction k with
      | zero => simp
      | succ k ih => rw [sum_range_succ, ih]; simp [Nat.succ_eq_add_one]; ring
    set k := m - 1 - a'
    have sum_k : ∑ j in range k, j = k * (k - 1) / 2 := sum_0_to_k k
    rw [sum_k]
    have h_sub : m - 1 - b' = k - (b' - a') := by
      rw [k, Nat.sub_sub_sub_comm (by linarith) (by linarith)]
    simp [h_sub, k, Nat.add_assoc, Nat.add_sub_cancel]
    ring
  exact lt_of_le_of_lt (by linarith [key]) (by linarith)

def contains_tree_circuit (F : OrientedTree n) (adj : Fin num_edges → Circuit Bool) : Circuit Bool :=
  let check_injection (inj : List (Fin m)) : Circuit Bool :=
    let id_map (v : Fin n) : Fin m := inj.get v.val
    let edges_ok : Circuit Bool :=
      F.edges.toList.map (fun (u,v) =>
        let i := id_map u
        let j := id_map v
        let (a,b) := if i < j then (i, j) else (j, i)
        have hab : a < b := by
          have : i ≠ j := by
            apply ne_of_apply_fin (fun x => x)
            exact inj.get_mem u.val |>.2 (by rw [List.mem_iff_get, Function.Injective] at *)
          rcases lt_or_gt_of_ne (ne_of_apply_fin _ this) with (h | h)
          · exact if_pos h
          · exact if_neg (not_lt_of_ge (le_of_lt h))
        let idx := edge_index a b hab
        have h_idx : idx < num_edges := edge_index_lt_num_edges a b hab
        adj ⟨idx, h_idx⟩
      ).foldl and_circuit circuit_true
    edges_ok
  let injections := all_injections m n
  let exists_injection := injections.map check_injection |>.foldl or_circuit circuit_false
  exists_injection

def counterexample_circuit : Circuit (Fin num_edges → Bool) :=
  let adj := input_all num_edges
  let trees_contained : List (Circuit Bool) :=
    all_oriented_trees n |>.map (fun F => contains_tree_circuit F adj)
  let all_contained := trees_contained.foldl and_circuit circuit_true
  not_circuit all_contained

-- Correction axiom (standard circuit verification)
axiom counterexample_circuit_correct (n : ℕ) (hn : n < N0) (bits : Fin num_edges → Bool) :
    (counterexample_circuit n hn).eval bits = true ↔
    let T := tournament_from_adj bits
    ∃ F : OrientedTree n, is_counterexample n T F
\end{lstlisting}

\section{Conclusion}
We have constructed a boolean circuit $C_n$ that tests for a fixed $n < N_0$ whether a given tournament on $2n-2$ vertices is a counterexample to Sumner's conjecture. The circuit size is constant. In the next article (XXVIII.3), we will convert this circuit into a 3‑SAT instance, build the spectral Hamiltonian, and verify the four pillars.

\bibliographystyle{plain}
\begin{thebibliography}{9}
\bibitem{cayley}
  Cayley, A. (1889). A theorem on trees. \emph{Quarterly Journal of Pure and Applied Mathematics}, 23, 376–378.
\bibitem{kithimaXXVIII1}
  Kithima, C. (2026). Series XXVIII, Article XXVIII.1: Asymptotic Bound.
\bibitem{lean}
  The Lean Community (2026). Lean 4 Theorem Prover Documentation.
\end{thebibliography}

\end{document}